% Bagian: computability
% Bab: computability-theory
% Subbagian: halting-problem

\documentclass[../../../include/open-logic-section]{subfiles}

\begin{document}

\olfileid[id]{cmp}{thy}{hlt}
\olsection{Masalah Penghentian}

Berdasarkan konstruksinya, fungsi dapat dihitung parsial universal
$\fn{Un}(e,x)$ terdefinisi jika dan hanya jika komputasi fungsi yang dikodekan
oleh~$e$ menghasilkan suatu nilai bagi masukan~$x$. Wajar untuk bertanya apakah
kita dapat memutuskan apakah hal ini terjadi. Ternyata kita tidak dapat
melakukannya. Dalam model komputasi mesin Turing, ini berarti pertanyaan apakah
suatu mesin Turing yang diberikan berhenti pada suatu masukan yang diberikan
tidak dapat diputuskan secara komputasional. Karena itu, teorema berikut dikenal
sebagai ``ketakterputusan masalah penghentian.'' Kita akan memberikan dua bukti
di bawah. Bukti pertama melanjutkan alur pembahasan sebelumnya, sedangkan bukti
kedua lebih langsung.

\begin{thm}
\ollabel{thm:halting-problem}
Misalkan
\[
h(e, x)  =
\begin{cases}
1 & \text{jika\/ $\fn{Un}(e, x)$ terdefinisi} \\
0 & \text{selain itu.}
\end{cases}
\]
Maka $h$ tidak dapat dihitung.
\end{thm}

\begin{proof}
Andaikan $h$ dapat dihitung. Kita menunjukkan bahwa hal ini memungkinkan kita
mendefinisikan fungsi dapat dihitung universal. Definisikan
\[
\fn{Un'}(e,x) =
\begin{cases}
\fn{Un}(e,x) & \text{jika $h(e,x) = 1$} \\
0 & \text{selain itu.}
\end{cases}
\]
Sekarang $\fn{Un'}(e, x)$ merupakan fungsi total dan dapat dihitung jika~$h$
dapat dihitung. Sebagai contoh, kita dapat mendefinisikan $g$ dengan rekursi
primitif sebagai berikut:
\begin{align*}
g(0, e, x) & \simeq 0 \\
g(y+1, e, x) & \simeq \fn{Un}(e,x);
\end{align*}
kemudian
\[
\fn{Un'}(e,x) \simeq g(h(e,x),e,x).
\]
Karena $\fn{Un'}(e,x)$ sama dengan $\fn{Un}(e,x)$ di setiap tempat fungsi
terakhir terdefinisi, $\fn{Un'}$ bersifat universal bagi fungsi-fungsi dapat
dihitung parsial yang kebetulan total. Namun, hal ini bertentangan dengan
\olref[nou]{thm:no-univ}.
\end{proof}

\begin{proof}
Andaikan $h(e,x)$ dapat dihitung. Definisikan fungsi $g$ dengan
\[
g(x) =
\begin{cases}
  0                & \text{jika $h(x,x) = 0$} \\
  \fundefined & \text{selain itu.}
\end{cases}
\]
Fungsi $g$ dapat dihitung secara parsial. Sebagai contoh, kita dapat
mendefinisikannya sebagai $\umin{y}{h(x,x) = 0}$. Jadi, untuk suatu~$e$,
$g(x) \simeq \fn{Un}(e, x)$ bagi setiap~$x$. Apakah $g$ terdefinisi pada $e$?
Jika ya, maka berdasarkan definisi~$g$, $h(e,e) = 0$ ($g$ hanya dapat mengambil
nilai~$0$ jika terdefinisi). Berdasarkan definisi~$h$, ini berarti bahwa
$\fn{Un}(e, e)$ tidak terdefinisi. Berdasarkan asumsi kita bahwa
$g(x) \simeq \fn{Un}(e, x)$ untuk setiap~$x$, kita memperoleh bahwa $g(e)$
tidak terdefinisi, suatu kontradiksi. Di lain pihak, jika $g(e)$ tidak
terdefinisi, maka $h(e,e) \neq 0$, sehingga $h(e,e) = 1$. Akibatnya,
$\fn{Un}(e, e)$ terdefinisi. Namun, karena $g(x) \simeq \fn{Un}(e, x)$, maka
$g(e)$ juga terdefinisi. Sekali lagi, kita memperoleh kontradiksi.
\end{proof}

\begin{tagblock}{TMs}
\begin{explain}
Kita dapat menggambarkan argumen ini dalam bahasa mesin Turing. Andaikan
terdapat suatu mesin Turing~$H$ yang menerima deskripsi mesin Turing~$E$ dan
masukan~$x$, lalu memutuskan apakah $E$ berhenti pada masukan~$x$. Kita lalu
dapat membangun mesin Turing lain~$G$ yang menerima satu masukan~$x$,
menjalankan $H$ untuk memutuskan apakah mesin $M_x$ dengan indeks~$x$ berhenti
pada masukan~$x$, dan melakukan hal sebaliknya. Dengan kata lain, jika $H$
melaporkan bahwa $M_x$ berhenti pada masukan~$x$, $G$ memasuki kalang tak
berhingga; jika $H$ melaporkan bahwa $M_x$ tidak berhenti pada masukan~$x$,
maka $G$ langsung berhenti. Apakah $G$ berhenti ketika indeksnya sendiri
diberikan sebagai masukan? Argumen di atas menunjukkan bahwa $G$ berhenti jika
dan hanya jika tidak berhenti---suatu kontradiksi. Jadi, anggapan kita bahwa
terdapat mesin Turing~$H$ semacam itu haruslah salah.
\end{explain}
\end{tagblock}

\end{document}
